\notepage{59}{Spin commutators and the spin-one-half matrices}
The full angular-momentum operator is
\[
 \mathbf J=\mathbf L\otimes I_s+I_r\otimes\mathbf S,
 \qquad \comm{J_i}{J_j}=\ii\hbar\sum_k\epsilon_{ijk}J_k.
\]
The abbreviated commutator $\comm{L_i}{S_j}$ means
\[
 \comm{L_i\otimes I_s}{I_r\otimes S_j}=0,
\]
because the two operators act on different factors. Expanding
$\comm{L_i+S_i}{L_j+S_j}$, the mixed terms vanish. Subtracting the
orbital commutator from the total one gives
\[
 \boxed{\comm{S_i}{S_j}=\ii\hbar\sum_k\epsilon_{ijk}S_k},\qquad
 \mathbf S\times\mathbf S=\ii\hbar\mathbf S.
\]

\smallheading{An electron has \texorpdfstring{$s=\tfrac12$}{s=1/2}}
The two $S_z$ eigenvalues are $+\hbar/2$ and $-\hbar/2$.
In the ordered basis
$\ket{+}=\ket{\tfrac12,\tfrac12}$,
$\ket{-}=\ket{\tfrac12,-\tfrac12}$,
\[
 S_z=\frac\hbar2\begin{bmatrix}1&0\\0&-1\end{bmatrix}.
\]
The ladder formula supplies the individual matrix entries,
\[
 \mel{\tfrac12,m'}{S_\pm}{\tfrac12,m}
 =\hbar\sqrt{(\tfrac12\mp m)(\tfrac12\pm m+1)}\,
   \delta_{m',m\pm1}.
\]
Consequently,
\[
 S_+=\hbar\begin{bmatrix}0&1\\0&0\end{bmatrix},\qquad
 S_-=\hbar\begin{bmatrix}0&0\\1&0\end{bmatrix},
\]
or equivalently
$S_\pm=(\hbar/2)\begin{bmatrix}0&1\pm1\\1\mp1&0\end{bmatrix}$.
It follows that
\[
 S_x=\tfrac12(S_++S_-)
 =\frac\hbar2\begin{bmatrix}0&1\\1&0\end{bmatrix},\qquad
 S_y=\tfrac1{2\ii}(S_+-S_-)
 =\frac\hbar2\begin{bmatrix}0&-\ii\\\ii&0\end{bmatrix}.
\]
The dimensionless matrices $\sigma_i=2S_i/\hbar$ are the Pauli
matrices; the $S_i$ themselves retain the units of angular momentum.
